Put \(p_{a,n}=q_{a,n}r_n\), \(\varepsilon=\varepsilon_{\mathrm{pil},n}\), and \(\bar D_n=2^{\bar J_y+d\bar J_z}\). Products preserve the declared anisotropic H\"older moduli, so the boundary-wavelet approximation bounds at the pilot resolutions are
\[
\max_a\|p_{a,n}-p_{a,n}^{\mathrm{proj}}\|_\infty
+\|r_n-r_n^{\mathrm{proj}}\|_\infty
\lesssim2^{-s_y\bar J_y}+2^{-s_z\bar J_z}\lesssim\varepsilon.   \tag{1}
\]
For every retained coefficient, expanding the exponential moment of its centered bounded empirical summands and using \(e^v-1-v\le v^2e^{|v|}/2\) gives Bernstein's inequality. A union bound over \(O(\{\bar D_n\})\) coefficients and finite wavelet overlap yield, uniformly over the model class,
\[
P\{E_n>C(\varepsilon+x)\}\le C\exp(-cnx^2/\bar D_n),
\quad 0\le x\le1,                                                \tag{2}
\]
where
\[
E_n=\max_e\left\{\|\widetilde r_n^{(-e)}-r_n\|_\infty
+\max_a\|\widetilde p_{a,n}^{(-e)}-p_{a,n}\|_\infty\right\}.
\]
The tuning identities in \(\mathrm{def{:}pilot\text{-}tuning\text{-}and\text{-}studentizer}\) give \(\bar D_n\asymp(n/\log n)^{1/(2\beta+1)}\). Thus (2) implies \(E_n=O_P(\varepsilon)\), and tail integration gives \(\|E_n\|_{L^k}\le C_k\varepsilon\) for each fixed \(k\).

On \(\widetilde r\ge c/2\),
\[
\frac{\widetilde p}{\widetilde r}-\frac pr
=\frac{\widetilde p-p}{\widetilde r}
+p\frac{r-\widetilde r}{r\widetilde r}.                          \tag{3}
\]
Equations (2)--(3) and the density lower bounds prove the conditional-pilot rate and \(P(\cap_e\mathcal E_n^{(-e)})\to1\). The construction normalizes every conditional density, including off this event. Since \(T(\widetilde r)\in[c/2,2C]\),
\[
\frac c{4C}\le\widehat r_n^{(-e)}\le\frac{4C}{c}.              \tag{4}
\]
Lebesgue orthonormality gives the deterministic Gram bounds. Moreover,
\[
\|\widehat G^{(-e)}-G\|_{\mathrm{op}}
\le\|\widehat r_n^{(-e)}-r_n\|_\infty,
\qquad
\widehat G^{-1}-G^{-1}=\widehat G^{-1}(G-\widehat G)G^{-1},      \tag{5}
\]
which proves the Gram perturbation assertions.

Let \(A_z=-\partial_y(\bar q_n\partial_y)\), \(A_z^{\mathrm{op}}=P_{J_y^{\mathrm{op}}}A_zP_{J_y^{\mathrm{op}}}\), and use hats for the pilot forms. Normalized clipping gives \(c/(4C)\le\widehat{\bar q}_n\le4C/c\). This, the true density bounds, and the zero-mean inequality \(\|v\|_2\le\|v'\|_2\) give uniform coercivity. The weak forms and (2) imply
\[
\|\widehat A_z^{\mathrm{op}}-A_z^{\mathrm{op}}\|_{H^1/\mathbb R\to H^{-1}}
\lesssim E_n.                                                     \tag{6}
\]
The resolvent identity gives an \(O_P(E_n+\lambda_n)\) perturbation in energy norm. If \(u=A_z^{-1}f\), integration of the Neumann equation gives \(u'=-F/\bar q_n\), where \(F(y)=\int_0^yf(v)dv\). Since \(\bar q_n^{-1}\) is uniformly \(a_y\)-H\"older, boundary-wavelet approximation and Galerkin orthogonality give, for retained \(f,g\),
\[
|\langle f,\{(A_z^{\mathrm{op}})^{-1}P_{J_y^{\mathrm{op}}}-A_z^{-1}\}g\rangle|
\lesssim2^{-2a_yJ_y^{\mathrm{op}}}\|f\|_2\|g\|_2
\le\varepsilon^{a_y}\|f\|_2\|g\|_2.                           \tag{7}
\]
The last inequality is the oversampling rule. Integrating against \(\widehat r\), using (2), and expanding both Gram inverses proves
\[
\|\widehat B^{(-e)}-B\|_{\mathrm{op}}
=O_P(\varepsilon^{a_y}+\varepsilon+\lambda_n).                   \tag{8}
\]

We now identify the term missed by the former expansion. Fix fold \(e\) and condition on its pilot sigma-field. Define the pilot joint law
\[
d\widehat P_{a,e}^{\dagger}(y,z)
=\widehat q_{a,n}^{(-e)}(y,z)\widehat r_n^{(-e)}(z)\,dy\,dz
\]
and the signed error \(\dot P_{a,e}=P_{a,n}-\widehat P_{a,e}^{\dagger}\); write \(d\widehat R_e(z)=\widehat r_n^{(-e)}(z)dz\) for their common pilot marginal. Both the true pair and pilot pair have a common covariate marginal.  Along the convex path \(P_{a,t}=\widehat P_{a,e}^{\dagger}+t\dot P_{a,e}\), that marginal remains common and all densities remain between fixed positive constants by (4) and \(\mathrm{ass{:}conditional\text{-}density\text{-}bounds}\). Hence the one-dimensional quantile formula is twice directionally differentiable along the path.

Let \(\widehat\Delta_e=\widehat q_{1,n}^{(-e)}-\widehat q_{0,n}^{(-e)}\), and let the actual raw discrepancy error in pilot conditional coordinates be
\[
\zeta_e(y,z)
=\frac{r_n(z)\Delta_n(y,z)-\widehat r_n^{(-e)}(z)\widehat\Delta_e(y,z)}
{\widehat r_n^{(-e)}(z)}.                                        \tag{9}
\]
It has zero outcome integral for every \(z\). If \(d_e\) denotes the vector of retained moments of \(\dot P_{1,e}-\dot P_{0,e}\), the Gram formula gives the exact identity
\[
d_e^\top\widehat B^{(-e)}d_e
=\widehat a_e^{\mathrm{op}}(\widehat\Pi_e\zeta_e,\widehat\Pi_e\zeta_e),       \tag{10}
\]
where \(\widehat a_e^{\mathrm{op}}(f,g)=\int f\widehat{\mathcal G}_{n,z}^{(-e)}g\,dy\,d\widehat R_e\) on the Galerkin space and \(\widehat\Pi_e\) is the \(L^2(dy\,d\widehat R_e)\) Gram projector. Direct expansion of the diagonal-free statistic also gives
\[
E(\widehat U_{e,J_y,J_z,n}\mid\mathscr F_{-e})
=d_e^\top\widehat B^{(-e)}d_e.                                  \tag{11}
\]

For clarity, we derive the comparison with the full Taylor term.  If \(F\) has density at least \(c_* >0\), a distribution-function perturbation of sup-norm \(E\) moves its inverse by at most \(E/c_*\). For \(a_y=\min\{s_y,1\}\),
\[
F(x+w)-F(x)-q(x)w
=\int_0^w\{q(x+t)-q(x)\}dt=O(|w|^{1+a_y}).                     \tag{12}
\]
Apply (12) to the two conditional quantile maps along \(P_{a,t}\), square their difference, and integrate first in the common marginal and then in \(t\). Exact Taylor's formula with integral remainder gives
\[
\theta_n=\widehat\theta_{\mathrm{plug}}^{(-e)}
+D\Theta_{\widehat P_e}[\dot P_e]
+2\int_0^1(1-t)\mathcal H_{e,t}(\dot P_e,\dot P_e)dt,           \tag{13}
\]
where \(\mathcal H_{e,t}\) is one half of the second directional derivative. The score in \(\mathrm{def{:}anisotropic\text{-}estimator}\) represents exactly the first derivative in (13), including the shared-marginal term. Let \(\widehat{\mathcal G}^{0}_{e,z}\) be the exact, unridged weighted Neumann inverse with coefficient \(\widehat{\bar q}_n^{(-e)}\), and put
\[
\widehat a_e^0(f,g)=\int f(y,z)\{\widehat{\mathcal G}^{0}_{e,z}g(\cdot,z)\}(y)dy\,d\widehat R_e(z).
\]
At \(t=0\), integration of the conditional Neumann equation identifies the equality Hessian with \(\widehat a_e^0\). Equations (3), (9), and (12) then give
\[
\left|2\int_0^1(1-t)\mathcal H_{e,t}(\dot P_e,\dot P_e)dt
-\widehat a_e^0(\zeta_e,\zeta_e)\right|
\le C\{E_n^{2+a_y}+\|\Delta_n\|_\infty^{a_y}E_n^2\}.          \tag{14}
\]
Indeed, the first term is the H\"older inverse-map remainder from (12); replacing the conditional tangent by the raw error (9), replacing the unequal-law Hessian by its equality inverse-elliptic form, and varying the common marginal each contribute at most the second term. Every division in this calculation is justified by the pathwise lower density bound.

The resolvent and oversampling bounds (6)--(8), applied to the retained coefficients of \(\zeta_e\), give
\[
\left|\widehat a_e^{\mathrm{op}}(\widehat\Pi_e\zeta_e,\widehat\Pi_e\zeta_e)
-\widehat a_e^0(\widehat\Pi_e\zeta_e,\widehat\Pi_e\zeta_e)\right|
\le C\{E_n^{2+a_y}+E_n^3\}.                                    \tag{15}
\]
Subtracting (13) from the conditional mean of the fold estimator and using (10)--(15) yields
\[
E(\widehat\theta_{e,J_y,J_z,n}-\theta_n\mid\mathscr F_{-e})
=-\{\widehat a_e^0(\zeta_e,\zeta_e)
-\widehat a_e^0(\widehat\Pi_e\zeta_e,\widehat\Pi_e\zeta_e)\}
+R_{0,e},                                                        \tag{16}
\]
with \(|R_{0,e}|\le C\{E_n^{2+a_y}+E_n^3+\|\Delta_n\|_\infty^{a_y}E_n^2\}\). Since \(a_y\le1\) and \(E_n\le1\) on the relevant event, \(E_n^3\le E_n^{2+a_y}\). Equation (16) is the required negative full-minus-projected sign. If the pilot is oracle, then \(\dot P_{a,e}=0\), hence \(\zeta_e=0\), and every term in (16) is zero.

It remains to bound the projection gap in (16). On the clipping-inactive event, the pilot wavelet estimates at their balanced resolutions have anisotropic H\"older norms bounded in every fixed \(L^k\): at the highest outcome level the stochastic sup-norm \(O_{L^k}(\varepsilon)\) is multiplied by \(2^{s_y\bar J_y}\asymp\varepsilon^{-1}\), and the covariate calculation is identical. Products, reciprocals of functions bounded below, and conditional normalization preserve these bounds. Thus
\[
\|\zeta_e\|_2=O_{L^k}(\varepsilon),
\qquad
\|\zeta_e\|_{\mathcal C^{s_y,s_z}}=O_{L^k}(1).                 \tag{17}
\]
The complement calculation in \(\mathrm{prop{:}anisotropic\text{-}spectrum\text{-}bias}\), now applied to \(\zeta_e\), gives
\[
|\widehat a_e^0(\zeta_e,\zeta_e)
-\widehat a_e^0(\widehat\Pi_e\zeta_e,\widehat\Pi_e\zeta_e)|
\le C\|\zeta_e\|_2
\{2^{-(s_y+1)J_y}+2^{-s_zJ_z}\}.                               \tag{18}
\]
For the quadratic complement terms, use both \(\|\zeta_e\|_2=O_{L^k}(\varepsilon)\) and the smooth approximation bounds, together with \(\min(x^2,y^2)\le xy\); this yields the factor \(\|\zeta_e\|_2\) in (18). The same inequality holds in \(L^2\); the exponentially small complement of the clipping-inactive event is harmless because all densities are clipped and all transport values are bounded by one. Equations (2), (14)--(18), and the anisotropic amplitude inequality
\[
\|\Delta_n\|_\infty
\lesssim\rho_n^{2/(2+3/s_y+d/s_z)}                              \tag{19}
\]
give the three deterministic terms in the claimed remainder bound. To verify (19), let \(A=\|\Delta_n\|_\infty\). The H\"older modulus keeps an amplitude proportional to \(A\) on an outcome interval of length proportional to \(A^{1/s_y}\) and a covariate rectangle of volume proportional to \(A^{d/s_z}\); zero outcome mass makes its antiderivative have size proportional to \(A^{1+1/s_y}\) on a comparable interval. The identity \(\rho_n^2=\int H_n^2/\bar q_n\) and the upper density bound then imply \(\rho_n^2\gtrsim A^{2+3/s_y+d/s_z}\).

Finally, subtract the conditional means in (16). The residual first-order coefficient has squared \(L^2(P_{0,n})\oplus L^2(P_{1,n})\) norm bounded by
\[
C\{E_n^{2+2a_y}+\|\Delta_n\|_\infty^{2a_y}E_n^2\}+o_P(\rho_n^2), \tag{20}
\]
by (12) and the energy-relative version of (8). Conditional independence makes its empirical average have variance \(m^{-1}\) times (20). Pair counting for the canonical kernel replacement gives
\[
E[\{U_e^{\mathrm{can}}(\widehat B)-U_e^{\mathrm{can}}(B)\}^2\mid\mathscr F_{-e}]
\le\frac C{m^2}\|\Gamma^{1/2}(\widehat B-B)\Gamma^{1/2}\|_{\mathrm{HS}}^2. \tag{21}
\]
At outcome level \(j\), (7)--(8) bound the Galerkin part by \(C\min\{2^{-2j},\varepsilon^{a_y}\}\), and therefore the ratio in (21) to \(\|K\|_{\mathrm{HS}}^2/m^2\) is \(o_P(1)\). Since \(\sqrt n\varepsilon\to\infty\), the square roots of (20)--(21) are absorbed by the deterministic terms already displayed plus
\[
o\!\left\{\frac{\rho_n}{\sqrt n}+\frac{A_{J_y,J_z,n}}n\right\}. \tag{22}
\]
Averaging the two folds proves the asserted \(L^2\) bound.

For the final covariance assertion, compact support gives \(\|\Phi(W)\|_2^2\lesssim D_{J_y,J_z}\). Exponential Markov bounds for scalar quadratic forms on a \(1/4\)-net yield
\[
\max_e\|\widehat\Gamma^{(-e)}-\Gamma\|_{\mathrm{op}}
=O_P\!\left(\sqrt{D_{J_y,J_z}\log n/m}+D_{J_y,J_z}\log n/m\right)=o_P(1). \tag{23}
\]
Uniform positive definiteness makes the square-root map locally Lipschitz. Equations (8), (21), and (23) prove the operator and relative Hilbert--Schmidt consistency of \(\widehat K\). Condition (23) is absent from (1)--(22), so it is needed only for an observed spectral plug-in.